\documentclass[11pt]{article}
\usepackage[margin=1in]{geometry}
\usepackage{amsmath,amssymb,amsthm,mathtools,bm}
\usepackage{enumitem}
\usepackage{booktabs}
\usepackage[hidelinks]{hyperref}

\newtheorem{theorem}{Theorem}
\newtheorem{proposition}[theorem]{Proposition}
\newtheorem{lemma}[theorem]{Lemma}
\newtheorem{corollary}[theorem]{Corollary}
\theoremstyle{definition}
\newtheorem{definition}[theorem]{Definition}
\newtheorem{assumption}{Assumption}
\renewcommand{\theassumption}{\Alph{assumption}}
\theoremstyle{remark}
\newtheorem{remark}[theorem]{Remark}

\newcommand{\R}{\mathbb{R}}
\newcommand{\N}{\mathbb{N}}
\newcommand{\E}{\mathbb{E}}
\newcommand{\cU}{\mathcal{U}}
\newcommand{\cV}{\mathcal{V}}
\newcommand{\cE}{\mathcal{E}}
\newcommand{\cG}{\mathcal{G}}
\newcommand{\cX}{\mathcal{X}}
\newcommand{\cO}{\mathcal{O}}
\newcommand{\cN}{\mathcal{N}}
\newcommand{\cF}{\mathcal{F}}
\newcommand{\ip}[2]{\langle #1,#2\rangle_W}
\newcommand{\nW}[1]{\lVert #1\rVert_W}
\newcommand{\n}[1]{\lVert #1\rVert}
\newcommand{\symW}{\operatorname{sym}_W}
\newcommand{\Ghat}{\widehat G}
\newcommand{\Lhat}{\widehat L}
\newcommand{\fhat}{\widehat f}
\newcommand{\Ehat}{\widehat E}

\title{Why regressing a policy onto measured-residual ILC steps converges\\ on the real system where model-based guided policy search does not}
\author{DistILC: formal analysis (draft, 2026-10-13)}
\date{}

\begin{document}
\maketitle

\begin{abstract}
We analyse the update at the core of DistILC: a goal-conditioned policy network is regressed, trial after trial, onto
one damped Gauss--Newton (Levenberg--Marquardt) iterative-learning-control (ILC) correction whose \emph{error} is
measured on the real system and whose \emph{Jacobian} comes from an approximate (nominal) model, evaluated along the
measured trajectory with the network's own feedback. We make every space and regularity assumption explicit and prove:
(i) every fixed point of the update is a zero of the \emph{real} error, whatever the size of the model error
(Proposition~\ref{prop:fixedpoints}); (ii) local linear convergence to such a zero, and almost-sure convergence under
zero-mean trial noise with Robbins--Monro steps, under a single \emph{positivity} condition coupling the true
sensitivity, the network's tangent kernel and the model's damped pseudo-inverse (Theorems~\ref{thm:local}
and~\ref{thm:stoch}); (iii) checkable sufficient conditions for positivity, including a small-gain bound and arbitrary
multiplicative mismatch (Lemma~\ref{lem:suff}). We contrast this with guided policy search (GPS): with the nominal
model, GPS's fixed points carry the model's zeroth-order residual as an irreducible real error
(Proposition~\ref{prop:gpsnominal}); with dynamics fitted from real samples, GPS is unbiased only once the local fits are
identifiable, which needs at least $n+m+1$ rollouts per condition and iteration (Proposition~\ref{prop:ident}); and with
a nominal prior in place of identifiable fits, any \emph{cost-side} regularization (control cost, entropy) turns the
approximate-gradient fixed point into a biased one, whereas \emph{step-side} damping does not
(Proposition~\ref{prop:regbias}). A first-order statement of cross-goal transfer answers when correcting the flown goals
corrects unflown ones (Proposition~\ref{prop:transfer}). Section~\ref{sec:critique} lists the criticisms we expect and
what the analysis does and does not answer.
\end{abstract}

\tableofcontents

%=====================================================================================================================
\section{Setting, spaces and standing assumptions}\label{sec:setting}

Throughout, all vector spaces are finite-dimensional real vector spaces; ``norm'' means the norm induced by the stated
inner product; linear maps between such spaces are identified with matrices in fixed bases; $\n{\cdot}$ on linear maps
is the induced operator norm. For a linear map $M$ on an inner-product space $(\cE,\ip{\cdot}{\cdot})$ we write
$\symW(M) := \tfrac12(M + M^{\star})$, with $M^\star = W^{-1}M^\top W$ the $W$-adjoint, and $\symW(M)\succeq \mu I$
means $\ip{e}{Me} \ge \mu\nW{e}^2$ for all $e\in\cE$.

\subsection{Spaces}
\begin{enumerate}[label=(S\arabic*),leftmargin=3em]
\item \textbf{Time.} A trial has a fixed finite horizon of $T\in\N$ control periods, $t=0,\dots,T-1$.
\item \textbf{State, action, observation, goal.} The state space $\cX\subseteq\R^{n}$ is open; the action space is
  $U=[-1,1]^m\subset\R^m$ (normalized, compact, convex); the observation space is $\cO\subseteq\R^{n_o}$ with a
  measurable observation map $h:\cX\times\cG\to\cO$, $C^1$ in $x$; the goal space $\cG\subset\R^{d}$ is compact. The
  method updates the policy at a finite set of \emph{trained goals} $\cG_{\rm tr}=\{g_1,\dots,g_N\}\subset\cG$.
\item \textbf{Trial input space.} $\cV:=\R^{Tm}$, sequences $v=(v_0,\dots,v_{T-1})$ of additive action perturbations,
  with the Euclidean inner product. Stacked over trained goals: $\cV^N$ with the product inner product.
\item \textbf{Error space.} Each period produces $q$ error rows $c(x_t,g)\in\R^q$ (\emph{task errors}: in the
  pogo-stick instance the ballistic apex predictor minus $h^\star$ and the speed minus $v^\star$, active in flight);
  $\cE:=\R^{Tq}$ with the weighted inner product $\ip{a}{b}=a^\top W b$, $W\succ 0$ (the inverse squared error scales).
  Stacked: $\cE^N$, block-diagonal weight. We let $\cE_g\subseteq\cE$ denote the subspace of rows that are active along
  the trial (masked rows are identically zero and are removed); all statements about $\cE$ refer to $\cE_g$.
\item \textbf{Parameter space.} $\Theta=\R^p$ with the Euclidean inner product; the policy is a map
  $\pi:\Theta\times\cO\times\cG\to\R^m$, $(\theta,o,g)\mapsto\pi_\theta(o,g)$, followed by the clip onto $U$.
\item \textbf{Probability.} A probability space $(\Omega,\cF,P)$ with a filtration $(\cF_k)_{k\ge0}$, $\cF_k$ the
  $\sigma$-algebra generated by all trials up to iteration $k$; per-trial noise $\xi_k$ takes values in a measurable
  space $(\Xi,\mathcal B)$; all maps composed with noise below are jointly measurable.
\end{enumerate}

\subsection{Systems}
\begin{enumerate}[label=(D\arabic*),leftmargin=3em]
\item \textbf{Real system.} $x_{t+1}=f(x_t,u_t)$ with $f:\cX\times U\to\cX$ (noise enters through measurement and initial
  state, see (N1)). It may be \emph{hybrid}: $f$ is the period map of a piecewise-smooth flow with guards
  (touchdown/liftoff).
\item \textbf{Nominal model.} $\fhat:\cX\times U\to\cX$, differentiable along the trajectories used (in the pogo
  instance: the same SLIP equations without leg damping, at a coarser step).
\item \textbf{Closed loop and the additive input.} For policy $\theta$, goal $g$ and $v\in\cV$, the trial from
  $x_0^g$ is $u_t=\operatorname{clip}_U(\pi_\theta(h(x_t,g),g)+v_t)$, $x_{t+1}=f(x_t,u_t)$. The trial map is
  \[E_g:\Theta\times\cV\to\cE,\qquad E_g(\theta,v)=(c(x_t,g))_{t=0}^{T-1},\qquad e_g(\theta):=E_g(\theta,0).\]
  $\Ehat_g$ is defined in the same way with $\fhat$.
\end{enumerate}

\subsection{Regularity}
\begin{assumption}[Smoothness in a neighbourhood]\label{as:smooth}
There is an open set $\cN\subseteq\Theta$ such that for every $g\in\cG_{\rm tr}$ and $\theta\in\cN$:
\begin{enumerate}[label=(\alph*)]
\item \emph{(interior actions)} the clip is inactive along the trial (actions in the interior of $U$);
\item \emph{(fixed, transversal mode sequence)} the hybrid trial visits the same sequence of modes for all
  $\theta\in\cN$, and every guard crossing is transversal (the guard's normal derivative along the flow is non-zero), so
  that event times are $C^1$ functions of the state by the implicit function theorem and the trial map is $C^1$
  (saltation matrices included in its derivative);
\item \emph{(Lipschitz derivative)} $\theta\mapsto D_\theta e_g(\theta)$ is Lipschitz on $\cN$ with constant $L_2$, and
  $\pi$ is $C^2$ in $(\theta,o)$ with bounded second derivatives on $\cN\times h(\cX,\cG)$.
\end{enumerate}
\end{assumption}
Assumption~\ref{as:smooth}(b) is the hybrid-systems price: across a change of the mode sequence (a missed touchdown, a
grazing contact) the trial map is not differentiable and no first-order argument applies; $\cN$ must lie inside one
mode-sequence cell. Saturation (a) is similar: on an active clip the sensitivity to that action is zero.

\subsection{Sensitivities}
\begin{definition}[True and nominal closed-loop sensitivities]\label{def:G}
Along the trial of $(\theta,g)$ with states $x_t$, actions $u_t$, let $A_t=\partial_x f$, $B_t=\partial_u f$ (with
saltation at events), $C_t=\partial_x c(x_t,g)$, and the policy feedback $K_t=\partial_o\pi_\theta\,\partial_x h$.
The true closed-loop sensitivity $G_g(\theta)=\partial_v E_g(\theta,0)\in L(\cV,\cE)$ has blocks
\[G[t,j]=C_t\,\Phi^{\rm cl}(t,j+1)\,B_j\ (j<t),\qquad G[t,j]=0\ (j\ge t),\qquad
\Phi^{\rm cl}(t,s)=\prod_{r=s}^{t-1}(A_r+B_rK_r).\]
The \emph{nominal} sensitivity $\Ghat_g(\theta)$ is the same formula with $\widehat A_t=\partial_x\fhat$,
$\widehat B_t=\partial_u\fhat$ evaluated \textbf{at the measured real states and actions} $(x_t,u_t)$, the same $C_t$ and
the same $K_t$.
\end{definition}
\begin{remark}
Because $\Ghat$ is evaluated along the measured trajectory with the network's own measured feedback, $G-\Ghat$ contains
\emph{only} the derivative mismatch $(\partial f-\partial\fhat)$ at the same points (and, for hybrid systems, the
approximation of saltation terms); it contains no state-mismatch term. The zeroth-order mismatch $f-\fhat$ never enters
the update: the error is measured.
\end{remark}

\begin{lemma}[Chain rule through the closed loop]\label{lem:chain}
Under Assumption~\ref{as:smooth}, define $J_g(\theta)\in L(\Theta,\cV)$ by $(J_g\delta\theta)_t=\partial_\theta
\pi_\theta(o_t,g)\,\delta\theta$ along the trial. Then $D_\theta e_g(\theta)=G_g(\theta)J_g(\theta)$.
\end{lemma}
\begin{proof}
A parameter change $\delta\theta$ changes the action at period $t$ by $\partial_\theta\pi\,\delta\theta+K_t\delta x_t$;
the first term is exactly an additive input $v_t=(J_g\delta\theta)_t$ applied to the closed loop with feedback $K_t$,
and the second is the closed loop's own response. Linearizing the recursion $\delta x_{t+1}=(A_t+B_tK_t)\delta
x_t+B_t v_t$, $\delta x_0=0$, and $\delta e_t=C_t\delta x_t$ gives $\delta e=G_gJ_g\delta\theta$.
Differentiability of each map is Assumption~\ref{as:smooth}.
\end{proof}

%=====================================================================================================================
\section{The DistILC update}\label{sec:update}

\paragraph{ILC step (per trained goal).} With damping $\lambda>0$, gain $\beta\in(0,1]$:
\begin{equation}\label{eq:ilc}
\Delta v_g=-\beta\,\Lhat_g\,e_g,\qquad \Lhat_g:=(\Ghat_g^\top W\Ghat_g+\lambda I)^{-1}\Ghat_g^\top W\ \in L(\cE,\cV).
\end{equation}
The damping is on the \emph{step} (Levenberg--Marquardt), not in the objective.

\paragraph{Regression (one network for all goals).} The targets are $y_{g,t}=\pi_{\theta_k}(o_{g,t},g)+\Delta
v_{g,t}$ at the measured observations. The network is refit, $\theta_{k+1}\approx\arg\min_\theta\sum_{g,t}\n{\pi_\theta
(o_{g,t},g)-y_{g,t}}^2+\rho\n{\theta-\theta_k}^2$, with $\rho\ge0$ a proximal weight (trust region, finite fitting).
Stacking $J=(J_{g_1};\dots;J_{g_N})\in L(\Theta,\cV^N)$ and $\Delta v=(\Delta v_{g_i})_i$, the first-order solution of
this fit is the kernel (minimum-norm, ridge) step
\begin{equation}\label{eq:reg}
\theta_{k+1}-\theta_k=J^\top(JJ^\top+\rho I)^{-1}\Delta v+r_k,
\end{equation}
where $r_k$ collects the inexactness of the actual optimizer (finite Adam steps, nonlinearity of $\pi$ in $\theta$,
replayed targets) and $\Theta_k:=JJ^\top\in L(\cV^N,\cV^N)$ is the empirical neural tangent kernel on the visited
(observation, goal) pairs. Define the \emph{realized-correction operator}
\[S_k:=\Theta_k(\Theta_k+\rho I)^{-1}\in L(\cV^N,\cV^N),\qquad 0\preceq S_k\preceq I,\]
whose off-diagonal blocks $S_{g_ig_j}$ couple goals: the correction asked at $g_j$ moves the policy at $g_i$.

\paragraph{Closed-loop error recursion.} By Lemma~\ref{lem:chain} and Taylor's theorem,
\begin{equation}\label{eq:recursion}
e_{k+1}=e_k-\beta\,M_k\,e_k+G_kJ_kr_k+R_k,\qquad M_k:=G_kS_k\Lhat_k\in L(\cE^N,\cE^N),
\end{equation}
with $G_k=\operatorname{blkdiag}(G_{g_i}(\theta_k))$, $\Lhat_k=\operatorname{blkdiag}(\Lhat_{g_i})$ and the Taylor
remainder $\nW{R_k}\le \tfrac{L_2}2\n{\theta_{k+1}-\theta_k}^2$.

\paragraph{From the implementation to the analysed update.} The implementation (\texttt{hw\_stage.py}) differs from
\eqref{eq:ilc}--\eqref{eq:reg} in four ways, each mapped onto a term of the analysis.
(i) \emph{rms cap}: the step is scaled down when its rms exceeds a cap, i.e.\ a state-dependent gain
$\beta_k\in(0,\beta]$; every statement below holds for any gain sequence in $(0,\mu/\Gamma^2]$.
(ii) \emph{replayed targets}: the network is fit to the targets of all trials so far with weights halving per iteration
of age; the difference from fitting the current targets only is part of the regression error $r_k$ (the analysed
update is the no-replay limit, \texttt{--age-decay 0}, which performed the same in the pogo sweep).
(iii) \emph{finite Adam fitting} instead of the kernel step: also $r_k$.
(iv) \emph{rollback}: a trial worse than the last kept one restores the previous parameters and halves the cap; this
is a safeguard outside the analysis (a projection onto a sublevel set of the measured cost), used in
Theorem~\ref{thm:stoch}'s localization remark.
The sim stage (Section~\ref{sec:sim}) uses the same update with $f=\fhat$ and many lanes per iteration.

%=====================================================================================================================
\section{Fixed points are zeros of the real error}

\begin{assumption}[Realizability and rank]\label{as:rank}
For $\theta\in\cN$: (a) $J(\theta)$ is surjective onto $\cV^N$ (the network can realize any correction at the visited
inputs; $\Theta_k\succ0$); (b) each $\Ghat_g$ has full row rank on $\cE_g$, i.e.\ $\dim\cE_g\le\dim\cV$ and the
model can, to first order, move every active error row.
\end{assumption}

\begin{proposition}[No model bias at fixed points]\label{prop:fixedpoints}
Under Assumptions~\ref{as:smooth}--\ref{as:rank}, with exact regression ($r_k=0$): $\theta$ is a fixed point of
\eqref{eq:ilc}--\eqref{eq:reg} if and only if $e_g(\theta)=0$ for every trained goal $g$. This holds for every nominal
model $\fhat$ satisfying Assumption~\ref{as:rank}(b), regardless of $\n{f-\fhat}$ or $\n{G-\Ghat}$.
\end{proposition}
\begin{proof}
$\theta_{k+1}=\theta_k$ iff $J^\top(\Theta+\rho I)^{-1}\Delta v=0$. Since $J$ is surjective, $J^\top$ is injective,
and $(\Theta+\rho I)^{-1}$ is invertible, so this holds iff $\Delta v=0$, i.e.\ $\Lhat_ge_g=0$ for all $g$. $\Lhat_g=
(\Ghat_g^\top W\Ghat_g+\lambda I)^{-1}\Ghat_g^\top W$ is injective on $\cE_g$ because $W\succ0$ and $\Ghat_g^\top$ is
injective (full row rank of $\Ghat_g$). Hence $e_g=0$. Conversely $e_g=0$ gives $\Delta v=0$.
\end{proof}
\begin{remark}
The statement is about the location of fixed points, not about reaching them; Section~\ref{sec:conv} gives the
convergence condition. The essential mechanism: the update is linear in a \emph{measured} residual, so the residual,
not a model prediction of it, must vanish at rest. Step-side damping $\lambda$ changes the path, not the fixed points.
\end{remark}

%=====================================================================================================================
\section{Convergence}\label{sec:conv}

\begin{assumption}[Positivity]\label{as:pos}
There are $\mu>0$ and $\Gamma<\infty$ such that for all $\theta\in\cN$: $\symW(M(\theta))\succeq\mu I$ on $\cE^N$ and
$\n{M(\theta)}_W\le\Gamma$, where $M=GS\Lhat$ as in \eqref{eq:recursion}.
\end{assumption}
Positivity says that the correction the nominal model asks for, after passing through the network's realization
operator $S$ and the \emph{true} closed loop $G$, reduces the measured error to first order, uniformly, and including the
interference between goals (the off-diagonal blocks of $S$).

\begin{lemma}[Sufficient conditions for positivity]\label{lem:suff}
Let $\Ghat_g^+:=\Ghat_g^\top(\Ghat_g\Ghat_g^\top)^{-1}$ be the right pseudo-inverse (full row rank,
Assumption~\ref{as:rank}(b)). With $A:=W^{1/2}\Ghat_g$, the push-through identity gives
$\Lhat_g=A^\top(AA^\top+\lambda I)^{-1}W^{1/2}$, hence $\Lhat_g\to A^\top(AA^\top)^{-1}W^{1/2}=\Ghat_g^+$ as
$\lambda\to0$ (the weight cancels in the limit) and $\Ghat_g\Ghat_g^+=I$ on $\cE_g$. Write $\Delta_g:=G_g-\Ghat_g$.
\begin{enumerate}[label=(\roman*)]
\item \emph{(Small gain.)} If $S=I$ (exact realization without cross-goal interference), $\lambda\to0$, and
  $\n{\Delta_g\Ghat_g^+}_W\le1-\mu'$ for all $g$, then $\symW(M)\succeq\mu' I$.
\item \emph{(Multiplicative mismatch.)} If $S=I$, $\lambda\to0$ and $G_g=P_g\Ghat_g$ with $\symW(P_g)\succeq\mu I$
  ($P_g$ any linear map on $\cE_g$ with positive-definite symmetric part, e.g.\ gain, mass or actuator-scale errors of
  any size), then $\symW(M)\succeq\mu I$.
\item \emph{(Interference.)} For general $S$, write $M=D+O$ with $D$ the block-diagonal and $O$ the off-diagonal part.
  If $\symW(D)\succeq\mu_D I$ and $\n{O}_W\le\omega<\mu_D$, then $\symW(M)\succeq(\mu_D-\omega)I$.
\item \emph{(Continuity in $\lambda$.)} If (i) or (ii) holds with margin $\mu'$, it holds with margin
  $\mu'-O(\lambda/\sigma_{\min}(\Ghat)^2)$ for $\lambda>0$.
\end{enumerate}
\end{lemma}
\begin{proof}
(i) $G\Ghat^+=(\Ghat+\Delta)\Ghat^+=I+\Delta\Ghat^+$ because $\Ghat\Ghat^+=I$ on $\cE_g$ (full row rank); then
$\ip{e}{(I+\Delta\Ghat^+)e}\ge\nW{e}^2-\n{\Delta\Ghat^+}_W\nW{e}^2\ge\mu'\nW{e}^2$.
(ii) $G\Ghat^+=P\Ghat\Ghat^+=P$. (iii) $\ip{e}{Me}\ge\ip{e}{De}-|\ip{e}{Oe}|\ge(\mu_D-\omega)\nW{e}^2$.
(iv) $\Lhat_g-\Ghat_g^+=-\lambda A^\top(AA^\top+\lambda I)^{-1}(AA^\top)^{-1}W^{1/2}$, whose norm is at most
$\lambda\n{A}\n{W^{1/2}}/\sigma_{\min}(A)^{4}$; compose with $\n{GS}$ bounded.
\end{proof}
\begin{remark}[Positivity is checkable]
Positivity concerns the map from the measured error to the realized error change; its Rayleigh quotients are observable
from consecutive trials: $\cos\angle(\Ghat\Delta v,\,e_{k+1}-e_k)$ (predicted vs measured change). The pogo-stick study
logs this: positive in 97\% of trial pairs for the nominal-sim network, median 0.88; across 125 robot runs the median
cosine predicts the hardware stage's gain (Spearman 0.63); the 180-degree control ($\Ghat\to-\Ghat$) gives 0\% success.
A sign-wrong model is exactly the case where Assumption~\ref{as:pos} fails, and the method fails, as predicted.
\end{remark}

\begin{theorem}[Local linear convergence]\label{thm:local}
Let Assumptions~\ref{as:smooth}--\ref{as:pos} hold; let $c_J:=\sup_{\cN}\n{J^\top(\Theta+\rho I)^{-1}}$,
$c_L:=\sup_\cN\n{\Lhat}$, $c_G:=\sup_{\cN}\n{GJ}$; suppose the regression error satisfies $\n{r_k}\le\varepsilon_k$.
Let $0<\beta\le\mu/\Gamma^2$ and $\delta>0$ with $L_2\,c_J^2c_L^2\,\beta\,\delta\le\mu/4$; let
$\bar\varepsilon:=\sup_k\varepsilon_k$ and $\varepsilon_\Sigma:=\sum_k\varepsilon_k<\infty$ with
$c_G\bar\varepsilon+L_2\bar\varepsilon^2\le\beta\mu\delta/4$. If $\theta_0\in\cN$, $\nW{e_0}\le\delta$, and the closed
ball of radius $R:=\tfrac{4c_Jc_L}{\mu}\big(\delta+c_G\varepsilon_\Sigma+L_2\bar\varepsilon\varepsilon_\Sigma\big)+
\varepsilon_\Sigma$ around $\theta_0$ lies in $\cN$, then
\[\nW{e_{k+1}}\le\Big(1-\frac{\beta\mu}{4}\Big)\nW{e_k}+c_G\varepsilon_k+L_2\varepsilon_k^2,\]
so $\nW{e_k}\to0$ linearly if $\varepsilon_k\equiv0$, and $\limsup_k\nW{e_k}\le\frac{4}{\beta\mu}\limsup_k
(c_G\varepsilon_k+L_2\varepsilon_k^2)$ in general.
\end{theorem}
\begin{proof}
From \eqref{eq:recursion}, $\nW{(I-\beta M)e}^2=\nW e^2-2\beta\ip{e}{Me}+\beta^2\nW{Me}^2\le(1-2\beta\mu+\beta^2
\Gamma^2)\nW e^2\le(1-\beta\mu)\nW e^2$ for $\beta\le\mu/\Gamma^2$, and $\sqrt{1-\beta\mu}\le1-\beta\mu/2$.
The step satisfies $\n{\theta_{k+1}-\theta_k}\le\beta c_Jc_L\nW{e_k}+\varepsilon_k$, so
$\nW{R_k}\le\tfrac{L_2}{2}(\beta c_Jc_L\nW{e_k}+\varepsilon_k)^2\le L_2\beta^2c_J^2c_L^2\nW{e_k}^2+L_2\varepsilon_k^2$.
With $\nW{e_k}\le\delta$ and the choice of $\delta$, $L_2\beta^2c_J^2c_L^2\nW{e_k}^2\le\tfrac{\beta\mu}{4}\nW{e_k}$.
Adding $\nW{GJr_k}\le c_G\varepsilon_k$ gives the recursion, valid as long as $\theta_k\in\cN$ and
$\nW{e_k}\le\delta$. Invariance, by induction: if $\nW{e_k}\le\delta$ then
$\nW{e_{k+1}}\le(1-\tfrac{\beta\mu}4)\delta+c_G\bar\varepsilon+L_2\bar\varepsilon^2\le\delta$. Summing the recursion,
$\sum_k\nW{e_k}\le\frac{4}{\beta\mu}\big(\nW{e_0}+c_G\varepsilon_\Sigma+L_2\bar\varepsilon\varepsilon_\Sigma\big)$, hence
$\sum_k\n{\theta_{k+1}-\theta_k}\le\beta c_Jc_L\sum_k\nW{e_k}+\varepsilon_\Sigma\le R$: every iterate lies in the ball,
inside $\cN$, where all constants are valid.
\end{proof}

\paragraph{Noise.}
\begin{enumerate}[label=(N\arabic*),leftmargin=3em]
\item The measured error is $\tilde e_k=e(\theta_k)+\xi_k$ with $\E[\xi_k\mid\cF_k]=0$ and
  $\E[\nW{\xi_k}^2\mid\cF_k]\le\sigma^2$ (trial-to-trial noise from initial states, disturbances and sensing).
\item The Jacobian used in iteration $k$ is $\cF_k$-measurable, or its dependence on $\xi_k$ is at most second order:
  $\nW{\E[(\Lhat(\tilde x_k)-\Lhat(x_k))\xi_k\mid\cF_k]}\le c_\xi\sigma^2$. (Our implementation evaluates $\Ghat$ at the
  measured states of the same trial; with idealized sensing in the pogo study $\xi\equiv$ start-state randomness only.)
\end{enumerate}

\begin{theorem}[Almost-sure convergence with Robbins--Monro gains]\label{thm:stoch}
Let Assumptions~\ref{as:smooth}--\ref{as:pos}, (N1), (N2) with $c_\xi=0$ hold, $r_k=0$, and use gains
$\beta_k>0$ with $\sum_k\beta_k=\infty$, $\sum_k\beta_k^2<\infty$, $\beta_k\le\mu/\Gamma^2$. On the event
$\Omega_\cN:=\{\theta_k\in\cN,\ \nW{e(\theta_k)}\le\delta\ \forall k\}$, $e(\theta_k)\to0$ almost surely. If
$c_\xi>0$, $\limsup_k\nW{e(\theta_k)}=O(c_\xi\sigma^2/\mu)$ a.s.\ on $\Omega_\cN$.
\end{theorem}
\begin{proof}
Let $V_k=\nW{e(\theta_k)}^2$. The update uses $-\beta_k M_k\tilde e_k=-\beta_kM_ke_k-\beta_kM_k\xi_k$; by (N1) and
$\cF_k$-measurability of $M_k$, $\E[\ip{e_k}{M_k\xi_k}\mid\cF_k]=0$. Expanding as in Theorem~\ref{thm:local},
$\E[V_{k+1}\mid\cF_k]\le(1-\beta_k\mu)V_k+C\beta_k^2(V_k+\sigma^2)$ on $\Omega_\cN$ for a constant $C$ collecting
$\Gamma$, $L_2$, $c_J$, $c_L$. The Robbins--Siegmund theorem gives $V_k\to V_\infty$ a.s.\ and
$\sum_k\beta_k\mu V_k<\infty$ a.s.; since $\sum\beta_k=\infty$, $V_\infty=0$. With $c_\xi>0$ the conditional mean of the
update carries an extra $O(\beta_kc_\xi\sigma^2)$ term, which leaves a neighbourhood of radius $O(c_\xi\sigma^2/\mu)$.
\end{proof}
\begin{remark}[Localization]
The event $\Omega_\cN$ is not automatic: a noisy step can leave $\cN$ (change the mode sequence, saturate). The
implementation's rollback (a trial worse than the last kept one restores the previous network and halves the step) is a
practical projection onto a safe set; it makes the iterate sequence a.s.\ remain in a sublevel set of the measured cost,
not in $\cN$. We state the theorem on $\Omega_\cN$ and do not claim more.
\end{remark}

%=====================================================================================================================
\section{Guided policy search: where it fails and where it does not}\label{sec:gps}

We take GPS in its standard form \cite{levine2013gps,levine2014unknown,montgomery2016mdgps}: per trained condition
(goal) a local controller $p_g$ is optimized by a trajectory optimizer (iLQR/LQG) on a cost $\ell$ under a dynamics
model, subject to a trust region (KL) to the current policy; the policy is regressed onto the local controllers; repeat.
Three cases cover the variants relevant to sim-to-real.

\subsection{GPS with the nominal model}
\begin{proposition}[Irreducible zeroth-order bias]\label{prop:gpsnominal}
Let GPS use $\fhat$ for its trajectory optimization, with the task cost $\ell=\tfrac12\nW{\Ehat_g}^2$ (any additional
cost terms only add bias). At any GPS fixed point $\theta_{\rm GPS}$ with exact regression, the local controllers are
model-optimal, $D_\theta\tfrac12\nW{\Ehat_g(\theta_{\rm GPS})}^2=0$; if the goals are feasible in the model and the
model-optimal set is $\{\Ehat_g=0\}$, then the real error is
\[e_g(\theta_{\rm GPS})=\underbrace{\Ehat_g(\theta_{\rm GPS})}_{=0}+\rho_g(\theta_{\rm GPS}),\qquad
\rho_g:=e_g-\Ehat_g,\]
with $\rho_g$ the zeroth-order sim-to-real residual along the policy,
so $\inf\{\nW{e_g(\theta)}:\Ehat_g(\theta)=0\}=\inf\{\nW{\rho_g(\theta)}:\Ehat_g(\theta)=0\}$, which is strictly positive
whenever the residual does not vanish on the model's solution set. No number of GPS iterations reduces it; DistILC's
fixed points (Proposition~\ref{prop:fixedpoints}) have zero real error.
\end{proposition}
\begin{proof}
Immediate from the definitions: the GPS update depends on the real system only through the regression's input
distribution, and its targets are model-optimal; the decomposition is the definition of $\rho_g$.
\end{proof}
\begin{remark}
This is the pogo study's zero-shot row: the nominal-sim network has 0\% success on the real sticks (apex 4.5 cm low:
the residual is the unmodelled leg damping), while the nominal-sim validation reports 100\%.
\end{remark}

\subsection{GPS with dynamics fitted from real samples}
In the unknown-dynamics variant \cite{levine2014unknown}, GPS fits a time-varying linear(-Gaussian) model
$x_{t+1}\approx F_t[x_t;u_t;1]$ per condition and iteration from $N_s$ real rollouts and optimizes in that model.
Because the fit includes the affine term from the real samples, the model reproduces the real trajectory to zeroth order,
and with accurate first-order fits the fixed points are real optima: \emph{this variant is unbiased}. Its cost is
first-order information:

\begin{proposition}[Identifiability needs $n+m+1$ rollouts per condition and iteration]\label{prop:ident}
For each $t$, the least-squares fit of $F_t\in\R^{n\times(n+m+1)}$ from samples
$\{(x_t^i,u_t^i,x_{t+1}^i)\}_{i=1}^{N_s}$ has a unique solution iff the design matrix $Z_t=[x_t^i;u_t^i;1]_{i}\in
\R^{(n+m+1)\times N_s}$ has full row rank, which requires $N_s\ge n+m+1$ (and persistently exciting, i.e.\ deliberately
perturbed, rollouts). For $N_s<n+m+1$, $F_t$ is determined on the $(n+m+1-\operatorname{rank}Z_t)$-dimensional
complement only by the prior/regularizer.
\end{proposition}
\begin{proof}
Normal equations $F_tZ_tZ_t^\top=X_{t+1}Z_t^\top$ have a unique solution iff $Z_tZ_t^\top$ is invertible iff
$\operatorname{rank}Z_t=n+m+1\le N_s$. For lower rank, any $F_t+\Delta$ with $\Delta Z_t=0$ fits equally well.
\end{proof}
With $N$ trained goals, unbiased first-order fits cost at least $N(n+m+1)$ real rollouts per iteration (the GMM prior of
\cite{levine2014unknown} pools across time and conditions and reduces this in practice, at the price of making the
first-order model prior-dominated). DistILC uses one rollout per goal per iteration and no fitted first-order model: it
takes the first-order information from the nominal model and needs only Assumption~\ref{as:pos} of it.

\subsection{GPS with a nominal prior and too few samples: the cost-regularization bias}
With $N_s$ below identifiability, a fitted model with a nominal prior has the measured zeroth order and an (essentially)
nominal first order: its trajectory optimization then computes an approximate-Jacobian step on the measured residual,
much like \eqref{eq:ilc}. The remaining difference is where regularization enters.
\begin{proposition}[Cost-side regularization biases approximate-Jacobian fixed points; step-side damping does not]\label{prop:regbias}
Let $\Phi(v)=\tfrac12\nW{E(v)}^2+\tfrac12 v^\top Rv$, $R\succeq0$ (control cost; entropy and KL-to-prior terms have
the same effect locally), and consider any method whose fixed points satisfy the approximate stationarity
$\Ghat^\top WE(\hat v)+R\hat v=0$ (e.g.\ iLQR/DDP run to convergence in a model with the measured zeroth order and the
nominal first order). Then at $\hat v$ the true gradient is
\[\nabla\Phi(\hat v)=(G-\Ghat)^\top W E(\hat v),\]
which is non-zero whenever $R\hat v\ne0$, $\Ghat$ has full row rank and $(G-\Ghat)^\top W$ is injective on the range
of $E(\hat v)$; so $\hat v$ is neither a real stationary point nor a real zero, with
$\n{\nabla\Phi(\hat v)}\le\n{G-\Ghat}_W\nW{E(\hat v)}$ and $\nW{E(\hat v)}\ge\n{R\hat v}/\n{\Ghat^\top W}$. With
$R=0$ and damping only in the step \eqref{eq:ilc}, the fixed points are $\{E=0\}$ (Proposition~\ref{prop:fixedpoints}),
the global minimizers of the real $\tfrac12\nW E^2$.
\end{proposition}
\begin{proof}
$\nabla\Phi=G^\top WE+Rv$; subtract the fixed-point condition. If $R\hat v\neq0$ then $\Ghat^\top WE(\hat v)=-R\hat
v\neq0$, so $E(\hat v)\ne0$ and $\nW{E}\ge\n{R\hat v}/\n{\Ghat^\top W}$.
\end{proof}
This is the precise sense of the contrast with Abbeel, Quigley and Ng \cite{abbeel2006inaccurate}: optimizing a utility
with approximate derivatives converges to a region of size proportional to the derivative error; root-finding a
measured residual with approximate derivatives converges to the root.

\subsection{Summary of the comparison}
\begin{center}
\begin{tabular}{@{}p{0.27\textwidth}p{0.15\textwidth}p{0.2\textwidth}p{0.3\textwidth}@{}}
\toprule
method & real rollouts per goal and iteration & first-order information & real error at fixed points\\
\midrule
GPS, nominal model & 0 & model & $\rho_g\neq0$ (Prop.~\ref{prop:gpsnominal})\\
GPS, fitted dynamics & $\ge n+m+1$ & fitted (unbiased) & $0$ (if feasible, identifiable)\\
GPS, nominal prior, $N_s<n+m+1$ & $N_s$ & $\approx$ model & biased if $R\ne0$ (Prop.~\ref{prop:regbias})\\
DistILC hardware stage & 1 & model, positivity only & $0$ (Prop.~\ref{prop:fixedpoints}, Thm.~\ref{thm:local})\\
\bottomrule
\end{tabular}
\end{center}
``GPS fails'' is therefore a statement about a regime: few real trials and a wrong model. With enough real samples to
identify local dynamics, GPS does not fail; it costs at least $n+m+1$ times more real rollouts per iteration.

%=====================================================================================================================
\section{Does correcting the flown goals correct the others?}\label{sec:transfer}

Let $g'\notin\cG_{\rm tr}$ and $J_{g'}$ its policy Jacobian along its own trial. The first-order change of its error
after one update is, by Lemma~\ref{lem:chain} and \eqref{eq:reg},
\begin{equation}\label{eq:transfer}
\Delta e_{g'}=-\beta\,G_{g'}\,S_{g',{\rm tr}}\,\Lhat\,e_{\rm tr}+O(\n{\Delta\theta}^2),\qquad
S_{g',{\rm tr}}:=J_{g'}J^\top(JJ^\top+\rho I)^{-1}.
\end{equation}

\begin{proposition}[First-order transfer]\label{prop:transfer}
Assume, for some $g\in\cG_{\rm tr}$ at goal distance $d=\n{g'-g}$, within one mode-sequence cell:
(T1) systematic error, $\nW{e_{g'}-e_g}\le L_ed$; (T2) sensitivity continuity, $\n{G_{g'}-G_g}\le L_Gd$; (T3) kernel
coupling, $S_{g',{\rm tr}}=\kappa\,S_{g,{\rm tr}}+\Sigma$ with $\kappa>0$ and $\n{\Sigma}\le s(d)$, $s(0)=0$;
(T4) positivity at $g$ with margin $\mu_g$ for the row of $M$ belonging to $g$. Then
\[-\ip{e_{g'}}{\Delta e_{g'}}\ \ge\ \beta\Big(\kappa\mu_g\nW{e_g}^2-c\,(L_ed+L_Gd+s(d))\,\nW{e_{\rm tr}}\,\nW{e_g}\Big)
-O(\beta^2),\]
so the unflown goal's error decreases to first order for $d<r^\star\propto\kappa\mu_g\nW{e_g}/(L_e+L_G+s'(0))\,
\nW{e_{\rm tr}}^{-1}$: transfer is local, with a radius set by how fast the error, the sensitivity and the network's
coupling vary with the goal.
\end{proposition}
\begin{proof}
Substitute (T1)--(T3) into the inner product
\[\ip{e_{g'}}{G_{g'}S_{g',{\rm tr}}\Lhat e_{\rm tr}},\]
expand around $\ip{e_g}{G_gS_{g,{\rm tr}}\Lhat e_{\rm tr}}\ge\mu_g\nW{e_g}^2$ (T4), and bound each cross term by
Cauchy--Schwarz with the uniform bounds $c_G,c_L$; the constant $c$ collects them.
\end{proof}
\paragraph{Evidence.} Pogo stick, hardware stage on one goal $(0.8,0.65)$ for 6 trials per robot, evaluated on a
$7\times5$ goal grid over the whole box: the error ratio after/before is $0.11$ at the trained goal, $0.16$--$0.45$
within normalized distance $0.3$, $0.37$--$0.89$ at $0.5$--$0.73$; the error decreased at all 35 goals. The main
experiment's 8 test goals are all unflown; all improved. The sources of a systematic error (T1) are robot-wide
(damping, mass, actuator scale), which is why transfer is broad here. Mode boundaries (e.g.\ a different contact
sequence) violate (T2)/(T3) and give no first-order guarantee; nor does the theory guarantee anything beyond the trained
goals: generalization is a property of the network and the data, not of the update.

%=====================================================================================================================
\section{The simulation stage}\label{sec:sim}
In simulation the ``real'' system is the nominal model, $G=\Ghat$ up to the hybrid approximation of saltation terms, and
the recursion \eqref{eq:recursion} has $M=GS\Lhat$ with $G\Lhat=G(G^\top WG+\lambda I)^{-1}G^\top W\succeq0$ (a
damped projector). Per lane positivity therefore holds automatically in the limit $S\to I$; what can fail is the
\emph{stacked} condition through $S$: lanes compete for one network (Lemma~\ref{lem:suff}(iii)), and regression error
$r_k$ accumulates. We observed this: the nominal-sim stage's speed loop drifts late (the nominal-sim score cannot see
it; a sim-only DR-family score can), and under domain randomization the drift starts before the DR objective is learned.
\begin{proposition}[Monotone acceptance]\label{prop:accept}
If each update is kept only when the measured cost $\sum_{\rm lanes}\tfrac12\nW{e}^2$ of the same lanes (same initial
states, models and exploration noise) decreases, the cost of the accepted sequence on the lanes it is evaluated on is
non-increasing, and an update is accepted whenever the first-order decrease
$\beta\mu\sum\nW{e}^2$ exceeds the second-order and regression terms of \eqref{eq:recursion}, i.e.\ for
$\beta$ small enough when Assumption~\ref{as:pos} holds on those lanes.
\end{proposition}
\begin{proof}
The acceptance test is exactly the sign of the realized change of the cost; the rejection branch restores $\theta_k$.
The second statement is Theorem~\ref{thm:local}'s one-step inequality summed over lanes.
\end{proof}
This is a guarantee on the evaluated lanes only (a fresh batch each iteration), not on the population; it removed the
drift in the pogo study (speed error $0.04$--$0.10$ m/s at iteration 60 against $0.7$ m/s without it).

%=====================================================================================================================
\section{Domain randomization, for comparison}
Let $\xi\sim\nu$ index the randomized family and $L(\theta;\xi)$ the task loss. DR's SGD converges (under the usual
conditions) to stationary points of $\bar L(\theta)=\E_\nu L(\theta;\xi)$. If reality is $\xi^\star$, the real loss at
DR's optimum is $L(\theta_{\rm DR};\xi^\star)$, which is small only if (C1) $\xi^\star$ lies where $\nu$ has mass and
(C2) the compromise optimum $\arg\min\bar L$ is near $\arg\min L(\cdot;\xi^\star)$ (no conflicting family members, or a
policy that can identify $\xi$ from observations). Neither (C1) nor (C2) is checkable from simulation alone; positivity
is checkable from a few real trial pairs. The two are not nested: reality can lie outside every family while positivity
holds (Lemma~\ref{lem:suff}(ii) with a large gain error), and a wide family can contain reality while the nominal
Jacobian points more than $90^\circ$ off. They are complementary: DR is a choice of warm start, the ILC update is a
convergence mechanism; the pogo study's best arm is the DR policy followed by the ILC hardware stage (94--96\%).

%=====================================================================================================================
\section{Anticipated criticisms}\label{sec:critique}
\begin{enumerate}[leftmargin=1.5em]
\item \textbf{``This is ILC / Abbeel et al.\ 2006 / GPS with a prior.''} The per-goal step \eqref{eq:ilc} is
  approximate-model ILC; we do not claim it. What is analysed here is its use as the regression target of one
  goal-conditioned feedback network shared across goals (operator $S$, Sections~\ref{sec:update}, \ref{sec:transfer}),
  the role of the network's own feedback in the sensitivity (Definition~\ref{def:G}), and the fixed-point contrast with
  cost-optimizing methods (Proposition~\ref{prop:regbias}). With a nominal prior and too few samples, GPS becomes a
  close relative; the difference is then the regularization (cost-side vs step-side) and the use of one step per trial
  instead of a converged local optimization.
\item \textbf{``Positivity is an unverifiable assumption.''} It is a condition on observables: predicted vs measured
  error changes over consecutive trials. Lemma~\ref{lem:suff} gives model-error conditions (small gain; multiplicative
  mismatch of any size). Empirically it holds in 97\% of pogo trial pairs and predicts the gain; the flipped-Jacobian
  control fails as predicted.
\item \textbf{``Only local.''} Yes: Theorems~\ref{thm:local}--\ref{thm:stoch} are local, on a neighbourhood $\cN$ inside one
  mode-sequence cell. A warm start in the basin (the simulation stage, or a DR policy) is part of the method.
\item \textbf{``Hybrid systems are not differentiable.''} Assumption~\ref{as:smooth}(b) restricts the analysis to a fixed
  transversal mode sequence, where the trial map is $C^1$ with saltation terms. The implementation's straight-through
  saltation terms are an approximation of $G$, i.e.\ part of $G-\Ghat$, covered by positivity; mode-sequence changes
  are outside the theory.
\item \textbf{``The network cannot realize the targets.''} Captured by $S\ne I$ (cross-goal coupling) and the regression
  error $r_k$ (Theorem~\ref{thm:local}'s $O(\varepsilon)$ neighbourhood). Realizability (Assumption~\ref{as:rank}(a)) is an
  over-parameterization assumption on the visited inputs; it can fail with conflicting targets (one observation, two
  goals asking for opposite corrections), which is the interference term of Lemma~\ref{lem:suff}(iii).
\item \textbf{``Guarantees only on the trained goals.''} Correct. Proposition~\ref{prop:transfer} is a first-order,
  conditional statement; generalization to unflown goals is empirical (Section~\ref{sec:transfer}).
\item \textbf{``What if a goal is infeasible on the real system?''} Then $\{e_g=0\}=\emptyset$ and
  Proposition~\ref{prop:fixedpoints} says there is no fixed point at that goal; any limit satisfies only
  $\Lhat e=0$ in the damped sense, which is \emph{not} the real least-squares stationarity $G^\top We=0$: our iteration
  is biased for infeasible goals, as Proposition~\ref{prop:regbias} shows for cost-optimizers. The implementation's stall
  guard and fallback ladders handle infeasible goals heuristically.
\item \textbf{``More error rows than inputs.''} Then $\Ghat$ cannot have full row rank and zero error is generically
  unreachable; the theory requires $\dim\cE_g\le\dim\cV$ (Assumption~\ref{as:rank}(b)). In the pogo instance $q=m=2$ and
  half the rows are masked in stance.
\item \textbf{``Noise in the Jacobian correlates with noise in the error.''} (N2) with $c_\xi>0$; the bias is second
  order in the noise level (Theorem~\ref{thm:stoch}). Averaging repeated trials or evaluating $\Ghat$ on the previous
  trial removes it.
\item \textbf{``Step sizes are tuned.''} The theory prescribes $\beta\le\mu/\Gamma^2$; the full Gauss--Newton step
  ($\beta=1$, small $\lambda$) was chosen on validation goals in the pogo study. Too large a step leaves the region of
  Theorem~\ref{thm:local}; the rollback guards against it.
\item \textbf{``The sensitivity uses the policy's own feedback, which changes.''} Yes: $M$ depends on $\theta$ through
  $K$ and $S$. A fragile network can lose positivity (observed: a DR network on the heavy pogo stick, median cosine
  $<0$, no gain). Assumption~\ref{as:pos} is uniform on $\cN$ for this reason.
\item \textbf{``The comparison to GPS is not empirical.''} Propositions~\ref{prop:gpsnominal}--\ref{prop:regbias} are
  structural. The nominal-model case is measured (zero-shot row); a fitted-dynamics GPS arm at equal real budget is the
  missing experiment.
\item \textbf{``DR wins zero-shot.''} When its family covers reality it can (PPO+DR 77\% vs 0\% for the nominal-sim
  network on the pogo sticks, whose damping lies inside the DR family). The claim is not dominance; it is that the ILC
  update converges on the real system under a different, checkable assumption, and that it composes with DR.
\item \textbf{``Finite horizon, finite trials.''} The theorems are asymptotic in the iteration count at fixed horizon $T$;
  at 24 trials the pogo study is in the transient regime, where the linear rate of Theorem~\ref{thm:local} is the
  relevant statement.
\end{enumerate}

%=====================================================================================================================
\section*{Assumption index}
\begin{center}
\begin{tabular}{@{}lp{0.72\textwidth}@{}}
\toprule
(S1)--(S6) & finite horizon; open state set, compact convex action set, compact goal set; trial input space $\R^{Tm}$
(Euclidean); error space $\R^{Tq}$ with $W\succ0$ (active rows only); parameter space $\R^p$; filtered probability
space, measurable noise.\\
(D1)--(D3) & real (possibly hybrid) dynamics, nominal model, closed-loop trial map with an additive input.\\
A & interior actions; fixed transversal mode sequence on $\cN$; Lipschitz $D_\theta e$; $C^2$ policy.\\
B & realizability ($J$ surjective on the visited inputs); full row rank of $\Ghat_g$ ($\dim\cE_g\le\dim\cV$).\\
C & positivity: $\symW(GS\Lhat)\succeq\mu I$, $\n{GS\Lhat}\le\Gamma$ on $\cN$, cross-goal interference included.\\
(N1)--(N2) & zero-mean, bounded-variance trial noise; Jacobian independent of the current noise (or second-order).\\
Thm.~\ref{thm:local} & small initial error, step $\beta\le\mu/\Gamma^2$, iterates stay in $\cN$, regression error $\varepsilon_k$.\\
Thm.~\ref{thm:stoch} & Robbins--Monro gains; on the event that the iterates stay in $\cN$.\\
Prop.~\ref{prop:transfer} & (T1)--(T4): systematic error, sensitivity continuity, kernel coupling, positivity at a
flown goal; same mode-sequence cell.\\
\bottomrule
\end{tabular}
\end{center}

\begin{thebibliography}{9}
\bibitem{levine2013gps} S.~Levine and V.~Koltun. Guided policy search. ICML, 2013.
\bibitem{levine2014unknown} S.~Levine and P.~Abbeel. Learning neural network policies with guided policy search under
unknown dynamics. NeurIPS, 2014. [verify]
\bibitem{montgomery2016mdgps} W.~Montgomery and S.~Levine. Guided policy search via approximate mirror descent.
NeurIPS, 2016. [verify]
\bibitem{abbeel2006inaccurate} P.~Abbeel, M.~Quigley and A.~Y.~Ng. Using inaccurate models in reinforcement learning.
ICML, 2006.
\bibitem{robbinssiegmund} H.~Robbins and D.~Siegmund. A convergence theorem for non negative almost supermartingales and
some applications. Optimizing Methods in Statistics, 1971.
\bibitem{kushner} H.~J.~Kushner and G.~G.~Yin. Stochastic Approximation and Recursive Algorithms and Applications.
Springer, 2003.
\end{thebibliography}
\end{document}
